\chapter{YAML and expressions}

YAML is not the hard part.
The hard part is how GitHub evaluates expressions and contexts.

\section{Quoting and multiline}
Use \code{|} for multiline shell scripts.
Prefer double quotes for strings with \code{:} or special characters.

\begin{lstlisting}[language=YAML]
- name: Multiline script
  run: |
    set -euo pipefail
    echo "Start"
    echo "Done"
\end{lstlisting}

\section{Contexts}
Contexts are objects available at runtime.
Common ones:
\code{github}, \code{env}, \code{inputs}, \code{secrets}, \code{runner}, \code{steps}, \code{jobs}.

\begin{lstlisting}[language=YAML]
- name: Print a few values
  run: |
    echo "ref=${{ github.ref }}"
    echo "sha=${{ github.sha }}"
\end{lstlisting}

\section{Expressions and conditionals}
Expressions live inside \code{\$\{\{ ... \}\}}.
You will use them in \code{if:}, \code{env:}, and \code{with:}.

\begin{lstlisting}[language=YAML]
- name: Only on main
  if: ${{ github.ref == 'refs/heads/main' }}
  run: echo "main branch"
\end{lstlisting}

\begin{warningbox}{Warning: strings vs booleans}
Inputs are strings by default.
If you need booleans, be explicit with comparisons like \code{inputs.flag == 'true'}.
\end{warningbox}

\section{Step outputs}
You can set outputs from a step and use them later.

\begin{lstlisting}[language=YAML]
- id: vars
  run: |
    echo "name=world" >> "$GITHUB_OUTPUT"

- name: Use output
  run: echo "Hello ${{ steps.vars.outputs.name }}"
\end{lstlisting}
